%lezione 19 del 4 maggio 2026 pag 42 appunti
%-----------------------------------------------------------------------------------------------------
%[ricorda giustificazione phi, giustifica calcolo flussi]
\subsection{Teorema ottico}
Si considera ora il flusso della corrente di probablità $\mathbf{J_\textbf{tot}}$: essa segue
l'equazione di continuità che, nel caso stazionario, è
$$\Phi_\Sigma(\mathbf{J_\textbf{tot}}) = \int_{\partial\Sigma}^{}d\Omega J^r_\text{tot}r^2 =0$$
dove $\Sigma$ è una sfera di raggio $r$ generico; il contributo di $\mathbf{J_\textbf{in}}$ vale
$$\Phi_\Sigma(\mathbf{J_\textbf{in}}) = \frac{\hbar}{m}r^2
\int_{0}^{2\pi}d\varphi\int_{-1}^{1}d\cos\vartheta k\cos\vartheta = 0$$
dunque si deve avere che
$$\int_{\partial\Sigma}^{}d\Omega \left(J^r_\text{out}+J^r_\text{interf}\right)r^2 = 0$$
Come già detto, la corrente di interferenza contribuisce solo per valori di $\vartheta\approx0$, se
ne calcola dunque il flusso attraverso una superficie definita da $\varphi\in[0,2\pi]$ e
$\vartheta\in[0,\delta\vartheta]$:
$$\Phi_\Sigma(\mathbf{J_\textbf{interf}})=
\int_{\partial\Sigma}d\Omega J^r_\text{interf}r^2 = \int_{0}^{2\pi}d\varphi
\int_{1-\frac12\delta\vartheta^2}^{1}d\cos\vartheta J^r_\text{interf}r^2$$
\textit{Osservazione} In questo calcolo la dipendenza da $\varphi$ di $f_k$ può essere
trascurata, in quanto, per valori piccoli di $\vartheta$, $f_k(\vartheta,\varphi)\approx
f_k(\vartheta,\varphi')$. Definendo $x:=\cos\vartheta$, il flusso vale
$$\Phi_\Sigma(\mathbf{J_\textbf{interf}}) = 
\frac{\hbar kr}{m} \int_{0}^{2\pi}d\varphi\Im\left[i\int_{1-\frac12\delta\vartheta^2}^{1}dxxf_k^*
e^{ikr(x-1)}+ f_ke^{ikr(1-x)}\right]
$$
poiché $x\approx1$, l'integrale in $dx$ è scrivibile come
$$f_k^*(0)I-f_k(0)I^*,\phantom{x}
I=\int_{1-\frac12\delta\vartheta^2}^{1}dx e^{ikr(x-1)} =
\frac{-i}{kr}\left(1-e^{-i\frac{kr}{2}\delta\vartheta^2}\right)\approx\frac{-i}{kr}$$
dove si è trascurato $e^{-i\frac{kr}{2}\delta\vartheta^2}$ per il solito lemma di Riemann.\\
Il flusso vale dunque
$$\Phi_\Sigma(\mathbf{J_\textbf{interf}})=
\frac{2\pi\hbar kr}{m}\Im\left[i\left(f^*_k(0)\frac{-i}{kr}+f_k(0)\frac{i}{kr}\right)\right]
=\frac{2\pi\hbar}{m}\Im\left[f^*_k(0) -f_k(0)\right]=-\frac{4\pi\hbar}{m}\Im(f_k(0))
$$
\textit{Osservazione} Per $\vartheta=0$, il valore di $\varphi$ non è definito, quindi si indica
$f_k(0)$ senza definire $\varphi$.

Il flusso di $\mathbf{J_\textbf{out}}$ vale infine
$$\Phi_\Sigma(\mathbf{J_\textbf{out}})=\frac{\hbar k}{m}\int_{\partial\Sigma}^{}d\Omega r^2
\frac{1}{r^2} \left|f_k(\vartheta,\varphi)\right|^2 =
\frac{\hbar k}{m}\int_{\partial\Sigma}^{}d\Omega \dv[]{\sigma}{\Omega} = \frac{\hbar k}{m}
\sigma_\text{tot}$$
dove $\sigma_\text{tot}$ è la sezione d'urto totale.

Imponendo l'equazione di continuità, si trova il cosiddetto \textbf{teorema ottico}:
$$\boxed{
\sigma_\text{tot} = \frac{4\pi}{k} \Im[f_k(0)]
}$$
che lega la sezione d'urto totale all'ampiezza differenziale.

Ricapitolando, si è ricavato che conoscere l'ampiezza differenziale significa conoscere sia la sezione
d'urto totale, sia quella differenziale.
\section{Approssimazione di Born}
È ora chiaro che lo studio del problema dello scattering elastico è equivalente alla determinazione
dell'ampiezza differenziale $f_k(\vartheta,\varphi)$: per capire come si può determinare questa
funzione, bisogna tornare all'equazione di Schrödinger iniziale.
\subsection{Equazione di Lippmann-Schwinger}
L'equazione di Schödinger del problema può essere scritta come
$$H\ket{E_+}=E\ket{E_+}$$
dove $H = H_0 + V = -\frac{\hbar^2}{2m}\laplacian + V$ e $\ket{E_+}$ è lo stato del sistema a energia
$E$. Questa si può riscrivere come
$$(H_0-E)\ket{E_+}=-V\ket{E_+} \rightarrow \ket{E_+} = (E-H_0)^{-1}V\ket{E_+} + \ket{E}$$
dove $\ket{E}$ è una soluzione dell'equazione libera con energia $E$.\\
L'inverso dell'operatore $E-H_0$ non è ben definito ovunque, in quanto esistono degli stati che
lo annullano (le particelle libere di energia $E$!); si definisce dunque\footnote{
    Ogni matematico rispettabile avrebbe un attacco cardiaco a vedere la definizione
    $A^{-1}=\frac{1}{A}$, ma ai fini di questo corso queste paranoie sono inutili.
}
\textit{l'operatore di Green}
$$G_+:=\lim\limits_{\varepsilon\to0}\frac{1}{E-H_0+i\varepsilon}$$
che è ovunque ben definito perché l'operatore $E-H_0$ è autoaggiunto e ha solo autovalori reali.

Con questo, si scrive ora \textbf{l'equazione di Lippmann-Schwinger}
$$\boxed{\ket{E_+} = \ket{E}+G_+V\ket{E_+}}$$
\textit{Osservazione} Questa è equivalente all'equazione integrale su cui si è lavorato in precedenza,
infatti, proiettandola nello spazio delle coordinate e inserendo due completezze, si ha
\begin{gather*}
\bra{\mathbf{x}}\ket{E_+} = \bra{\mathbf{x}}\ket{E} + \int_{}^{}d^3\mathbf{y}d^3\mathbf{z}
\bra{\mathbf{x}}G_+\ket{\mathbf{y}}\bra{\mathbf{y}}V\ket{\mathbf{z}}\bra{\mathbf{z}}\ket{E_+}
\rightarrow\\ \rightarrow\psi_{E_+}(\mathbf{x})=e^{i\mathbf{k}\cdot\mathbf{x}}+\int_{}^{}d^3\mathbf{y}
G_+(\mathbf{x},\mathbf{y})V(\mathbf{y})\psi_{E_+}(\mathbf{y})
\end{gather*}
che è l'equazione di prima, con l'identificazione $G_+(\mathbf{x},\mathbf{y})V(\mathbf{y}) =
G(\mathbf{x},\mathbf{y})U(\mathbf{y})$.

Invertendo l'equazione appena trovata, si ha
$$(\mathbb{I}-G_+V)\ket{E_+} = \ket{E}\rightarrow\ket{E_+}= \frac{1}{\mathbb{I}-G_+V}\ket{E}$$
ricordando che tutte le funzioni non polinomiali di operatori sono definite in serie, l'espressione
di sopra denota la serie geometrica: si trova quindi un'espressione \textit{perturbativa} per
$\ket{E_+}$
$$\boxed{\ket{E_+} = \sum_{n=0}^{\infty}(G_+V)^n\ket{E} = \ket{E} + G_+V\ket{E} +...}$$
\textbf{L'approssimazione di Born} consiste nel considerare solo i primi due termini di questa serie
perturbativa
$$\ket{E_+}_B = \ket{E}+G_+V\ket{E}$$
proiettando sullo spazio delle coordinate si trova
$$\psi_{E_+}(\mathbf{x}) = e^{i\mathbf{k}\cdot\mathbf{x}} - \frac{1}{4\pi}\int_{}^{}d^3\mathbf{y}
\frac{e^{i\mathbf{k}\cdot(\mathbf{x}-\mathbf{y})}}{|\mathbf{x}-\mathbf{y}|}U(\mathbf{y})
e^{i\mathbf{k}\cdot\mathbf{y}}$$
come in precedenza, se si considera $|\mathbf{x}|>>|\mathbf{y}|$ e $V$ con range finito si possono
approssimare $|\mathbf{x}-\mathbf{y}| = |\mathbf{x}|=r$ e $i\mathbf{k}\cdot(\mathbf{x}-\mathbf{y})=
ikr - ik\mathbf{\hat x}\cdot\mathbf{y}:=ikr-i\mathbf{k'}\cdot\mathbf{y}$ e si trova
\begin{gather*}
\psi_{E_+}(\mathbf{x}) = E^{i\mathbf{k}\cdot\mathbf{x}} + \frac{e^{ikr}}{r}f(\vartheta,\varphi),
\phantom{x}\text{con}\\f(\vartheta,\varphi)
=-\frac{1}{4\pi}\int_{}^{}d^3\mathbf{y}e^{-i\mathbf{k'}\cdot\mathbf{y}}U(\mathbf{y})e^{i\mathbf{k}
\cdot\mathbf{y}} = -\frac{m}{2\pi\hbar^2}\int_{}^{}d^3\mathbf{y}V(y)e^{-i(\mathbf{k'}-\mathbf{k})\cdot
\mathbf{y}} = -\frac{m\sqrt{2\pi}}{\hbar^2} \tilde V(\mathbf{k'}-\mathbf{k})
\end{gather*}
dove con $\tilde V$ si denota la \textit{trasformata di Fourier} rispetto all'impulso scambiato
$\mathbf{k'}-\mathbf{k}$ del potenziale di interazione.

\textit{Osservazione} L'ordine 0 dell'espressione perturbativa considera l'onda non interagente
($\ket{E}$), il primo ordine considera l'onda che ha interagito una sola volta con il potenziale e
il secondo considera anche l'onda che ha interagito due volte con il potenziale V e così via.
\subsection{Regime dell'approssimazione}
Come ogni approssimazione in fisica, anche quella di Born è sufficientemente precisa solo sotto certe
condizioni. Se la serie
$$\sum_{n=0}^{\infty}(G_+V)^n\ket{E}$$ fosse una serie numerica, essa convergerebbe se e solo se
$|G_+V|<1$: sulla falsa riga di questa condizione, si può dire che, qualitativamente, $G_+V$ dev
essere "piccolo"
$$G_+V = O\left(\frac{V}{E}\right) = O\left(\frac{V}{T+V}\right)<<1 \leftrightarrow T>>V$$
Più precisamente, si può definire la "grandezza" del potenziale $V$ guardando quanto esso distorca la
funzione d'onda: il fattore di distorsione della $\psi$ si definisce come
$$|c(k)|^2:=\left|\frac{\psi_\text{out}(0)}{\psi_\text{in}(0)}\right|^2 =|\psi_\text{out}(0)|$$
Nell'approssimazione di Born e con l'ulteriore ipotesi che $V$ sia centrale (come la maggior parte
dei potenziali di interesse in 3 dimensioni), $\psi_\text{out}(0)$ vale
\begin{gather*}
\psi_\text{out}(0) = -\frac{m}{2\pi\hbar^2}\int_{}^{}d^3\mathbf{y} \frac{e^{ik|\mathbf{y}|}}
{|\mathbf{y}|}e^{i\mathbf{k}\cdot\mathbf{y}}V(\mathbf{y}) = -\frac{m}{2\pi\hbar^2}
\int_{0}^{\infty}drr^2\int_{0}^{2\pi}d\varphi\int_{-1}^{1}d\cos\vartheta \frac{e^{ikr}}{r}
e^{ikr\cos\vartheta}V(r)=\\
=-\frac{m}{\hbar^2}\int_{0}^{\infty}drre^{ikr}V(r)\int_{-1}^{1}dxe^{ikrx}=
-\frac{m}{\hbar^2}\int_{0}^{\infty}drre^{ikr}V(r)\left(\frac{e^{ikr}-e^{-ikr}}{ikr}\right)=\\=
-\frac{2m}{\hbar^2k}\int_{0}^{\infty}e^{ikr}\sin(kr)V(r)
\end{gather*}
dunque
$$|c(k)|^2<<1\leftrightarrow|\psi_\text{out}(0)|\leq\frac{2m}{\hbar^2k}\int_{0}^{\infty}dr
\left|e^{ikr}\sin(kr)V(r)\right|\leq\frac{2m}{\hbar^2k}\int_{0}^{\infty}|V(r)|<
\frac{2m}{\hbar^2k}AV_0<<1
$$
dove $A$ è il range del potenziale e $V_0 = \underset{r<A}\max\,|V(r)|$. Questa condizione si può
riscrivere come
$$\frac{2m}{\hbar p}AV_0<<1\leftrightarrow p>>\frac{2mAV_0}{\hbar}$$
il principio di indeterminazione sancisce $pA\sim\hbar$, dunque
$$p>>\frac{2m}{p}V_0\leftrightarrow \frac{p^2}{2m}>>V_0$$
L'approssimazione di Born vale dunque per scattering ad alte energie cinetiche.
